% Part: first-order-logic
% Chapter: syntax-and-semantics
% Section: first-order-languages

\documentclass[../../../include/open-logic-section]{subfiles}

\begin{document}

% SEGMENT OLP-0151-S01
\olfileid{fol}{syn}{fol}

\olsection{முதல்தர மொழிகள்}

முதல்தர தருக்கத்தின் வெளிப்பாடுகள், \emph{!!{variable}s},
\emph{!!{constant}s}, \emph{!!{predicate}s}, சில வேளைகளில்
\emph{!!{function}s} ஆகியவற்றைக் கொண்ட அடிப்படைச் சொற்களஞ்சியத்திலிருந்து
உருவாக்கப்படுகின்றன. இவற்றுடன் தருக்கத் தொடுப்புக்குறிகள், அளவையடைகள்,
அடைப்புக்குறிகள் மற்றும் காற்புள்ளி போன்ற நிறுத்தக் குறிகளையும் சேர்த்து
\emph{சொற்களும்} \emph{!!{formula}s} உம் உருவாகின்றன.

\begin{explain}
முறைசாரா வகையில், !!{predicate}s பண்புகளுக்கும் தொடர்புகளுக்கும் பெயர்களாகவும்,
!!{constant}s தனிப்பட்ட பொருட்களுக்குப் பெயர்களாகவும், !!{function}s
வரைபடங்களுக்கு பெயர்களாகவும் உள்ளன. !!{identity}~$\eq$ ஐத் தவிர இவை
\emph{தருக்கமல்லாத குறிகள்}; இவை சேர்ந்து ஒரு மொழியை உருவாக்குகின்றன.
எந்த முதல்தர மொழி~$\Lang L$ உம் அதன் தருக்கமல்லாத குறிகளால் தீர்மானிக்கப்படுகிறது.
மிகப் பொதுவான நிலையில், $\Lang L$ ஒவ்வொரு வகையிலும் எண்ணற்ற குறிகளைக் கொண்டிருக்கலாம்.
\end{explain}

பொதுவான நிலையில், முதல்தர தருக்கத்தில் பின்வரும் குறிகளைப் பயன்படுத்துகிறோம்:

\begin{enumerate}
\item தருக்கக் குறிகள்
\begin{enumerate}
\item தருக்கத் தொடுப்புக்குறிகள்:
  \startycommalist
  \iftag{prvNot}{\ycomma $\lnot$ (மறுப்பு)}{}%
  \iftag{prvAnd}{\ycomma $\land$ (இணையல்)}{}%
  \iftag{prvOr}{\ycomma $\lor$ (பிரிப்பிணைவு)}{}%
  \iftag{prvIf}{\ycomma $\lif$ (!!{conditional})}{}%
  \iftag{prvIff}{\ycomma $\liff$ (!!{biconditional})}{}%
  \iftag{prvAll}{\ycomma $\lforall$ (பொதுஅளவையடை)}{}%
  \iftag{prvEx}{\ycomma $\lexists$ (இருப்பியல் அளவையடை)}{}.
\tagitem{prvFalse}{!!{falsity}~$\lfalse$ க்கான கூற்றுத் தருக்க மாறிலி.}{}
\tagitem{prvTrue}{!!{truth}~$\ltrue$ க்கான கூற்றுத் தருக்க மாறிலி.}{}
\item இரு-இட !!{identity}~$\eq$.
\item !!{variable}s இன் !!{denumerable} தொகுப்பு: $\Obj v_0$, $\Obj v_1$, $\Obj
  v_2$, \dots
\end{enumerate}
\item தருக்கமல்லாத குறிகள்; இவை முதல்தர தருக்கத்தின் \emph{நிலையான
  மொழியை} உருவாக்குகின்றன.
\begin{enumerate}
\item ஒவ்வொரு $n>0$ க்கும் $n$-இட !!{predicate}s இன் !!{denumerable} தொகுப்பு:
  $\Obj A^n_0$, $\Obj A^n_1$, $\Obj A^n_2$, \dots
\item !!{constant}s இன் !!{denumerable} தொகுப்பு: $\Obj c_0$, $\Obj c_1$, $\Obj
  c_2$, \dots.
\item ஒவ்வொரு $n>0$ க்கும் $n$-இட !!{function}s இன் !!{denumerable} தொகுப்பு:
  $\Obj f^n_0$, $\Obj f^n_1$, $\Obj f^n_2$, \dots
\end{enumerate}
\item நிறுத்தக் குறிகள்: (, ), மற்றும் காற்புள்ளி.
\end{enumerate}

நமது வரையறைகளும் முடிவுகளும் பெரும்பாலும் முதல்தர தருக்கத்தின் முழு நிலையான
மொழிக்காக வடிவமைக்கப்படும். இருப்பினும், பயன்பாட்டைப் பொறுத்து மொழியைச் சில
!!{predicate}s, !!{constant}s மற்றும் !!{function}s களுக்கு மட்டுமே கட்டுப்படுத்தலாம்.

\begin{ex}
கணிதத்தின்~$\Lang L_A$ மொழியில் ஒரே ஒரு இரு-இட !!{predicate}~$<$,
ஒரே ஒரு !!{constant}~$\Obj 0$, ஒரு ஓரிட !!{function}~$\prime$, மற்றும்
இரண்டு இரு-இட !!{function}s~$+$ மற்றும்~$\times$ உள்ளன.
\end{ex}

\begin{ex}
கணக்குத்தொகுப்புக் கோட்பாட்டின்~$\Lang L_Z$ மொழியில் ஒரே ஒரு இரு-இட
!!{predicate}~$\in$ மட்டுமே உள்ளது.
\end{ex}

\begin{ex}
வரிசைகளின்~$\Lang L_\le$ மொழியில் ஒரே ஒரு இரு-இட !!{predicate}~$\le$ மட்டுமே உள்ளது.
\end{ex}

மீண்டும், இவை மரபுகள் மட்டுமே: அதிகாரப்பூர்வமாக இவை மாற்றுப் பெயர்கள்.
உதாரணமாக, $<$, $\in$, $\le$ ஆகியவை முறையே $\Obj A^2_0$,
$\Obj 0$ என்பது $\Obj c_0$ க்கும், $\prime$ என்பது $\Obj f^1_0$ க்கும்,
$+$ என்பது $\Obj f^2_0$ க்கும், $\times$ என்பது $\Obj f^2_1$ க்கும் மாற்றுப் பெயர்கள்.

\iftag{defNot,defOr,defAnd,defIf,defIff,defTrue,defFalse,defEx,defAll}{%
முன்னர் அறிமுகப்படுத்திய முதன்மைத் தொடுப்புக்குறிகளுக்கும்
\iftag{notprvEx,notprvAll}{அளவையடைக்கும்}{அளவையடைகளுக்கும்} மேலாக,
பின்வரும் \emph{வரையறுக்கப்பட்ட} குறிகளையும் பயன்படுத்துகிறோம்:
\startycommalist
  \iftag{defNot}{\ycomma $\lnot$ (மறுப்பு)}{}%
  \iftag{defAnd}{\ycomma $\land$ (இணையல்)}{}%
  \iftag{defOr}{\ycomma $\lor$ (பிரிப்பிணைவு)}{}%
  \iftag{defIf}{\ycomma $\lif$ (!!{conditional})}{}%
  \iftag{defIff}{\ycomma $\liff$ (!!{biconditional})}{}%
  \iftag{defAll}{\ycomma $\lforall$ (பொதுஅளவையடை)}{}%
  \iftag{defEx}{\ycomma $\lexists$ (இருப்பியல் அளவையடை)}{}%
  \iftag{defFalse}{\ycomma !!{falsity}~$\lfalse$}{}%
  \iftag{defTrue}{\ycomma !!{truth}~$\ltrue$}}{}.

\begin{tagblock}{defNot,defOr,defAnd,defIf,defIff,defTrue,defFalse,defEx,defAll}
\begin{explain}
வரையறுக்கப்பட்ட குறி அதிகாரப்பூர்வமாக மொழியின் பகுதியல்ல; அது முறைசாரா
சுருக்கமாக அறிமுகப்படுத்தப்படுகிறது. முதன்மைக் குறிகளை மட்டும் பயன்படுத்தினால்
நீளமாகிவிடும் சூத்திரங்களைச் சுருக்க இது உதவுகிறது. அதைவிடப் பெரிய நன்மை,
நிறுவல்கள் சுருங்குகின்றன என்பதே. ஒரு குறி முதன்மையானதாக இருந்தால், நிறுவலில்
அதற்கெனத் தனி நிலை தேவைப்படும்; ஆகவே முதன்மைக் குறிகள் அதிகமானால் நிறுவல்களும்
நீளமாகும்.
\end{explain}
\end{tagblock}

% Alternate symbols

\begin{intro}
மேலே பயன்படுத்திய சொற்களும் குறிகளும் மாறுபட்டதாக உங்களுக்கு அறிமுகமாகி
இருக்கலாம். தருக்க நூல்களும் ஆசிரியர்களும் ``மறுப்பு''க்கு $\sim$, $\neg$,
அல்லது~! ஐயும், ``இணையல்''க்கு $\wedge$, $\cdot$, அல்லது~$\&$ ஐயும்
பொதுவாகப் பயன்படுத்துகின்றனர். ``நிபந்தனை'' அல்லது ``உட்கிடை''க்கான
பொதுக் குறிகள் $\rightarrow$, $\Rightarrow$, $\supset$ ஆகும்.
\iftag{prvIff,defIff}{``இருநிபந்தனை'', ``இரு-உட்கிடை'' அல்லது
  ``(பொருள்சார்) சமத்துவம்'' என்பவற்றுக்கு $\leftrightarrow$,
  $\Leftrightarrow$, $\equiv$ ஆகிய குறிகள் பயன்படுகின்றன.}{}
\iftag{prvFalse,defFalse}{$\lfalse$ குறி ``falsity'', ``falsum'',
  ``absurdity'' அல்லது ``bottom'' என்றும் அழைக்கப்படுகிறது.}{}
\iftag{prvTrue,defTrue}{$\ltrue$ குறி ``truth'', ``verum'' அல்லது
  ``top'' என்றும் அழைக்கப்படுகிறது.}{}

!!{constant}s (சில நேரம் பெயர்கள் என்றும் அழைக்கப்படுபவை) க்கு இலத்தீன்
எழுத்துமாலையின் தொடக்கத்திலுள்ள சிறிய எழுத்துக்களையும் (எ.கா. $a$, $b$, $c$),
!!{variable}s க்கு முடிவிலுள்ள சிறிய எழுத்துக்களையும் (எ.கா. $x$, $y$, $z$)
பயன்படுத்துவது மரபு. அளவையடைகள் !!{variable}s உடன் சேர்கின்றன; உதாரணமாக
$x$. பொதுஅளவையடைக்கு $\forall x$, $(\forall x)$, $(x)$, $\Pi x$,
$\bigwedge_x$ என்றும், இருப்பியல் அளவையடைக்கு $\exists x$, $(\exists
x)$, $(Ex)$, $\Sigma x$, $\bigvee_x$ என்றும் குறியீட்டு மாறுபாடுகள் உள்ளன.
\end{intro}

\begin{explain}
கூற்றுத் தருக்கத்தின் எல்லா இயக்கிகளையும் இரு அளவையடைகளையும் மொழியின்
முதன்மைக் குறிகளாகக் கருதலாம். அல்லது குறைந்த முதன்மைக் குறிகளின் தொகுப்பைத்
தேர்ந்தெடுத்து, பிற !!{operator}s ஐ வரையறுக்கப்பட்டவையாகக் கருதலாம்.
``மெய்மைச் சார்பில் முழுமையான'' பூலிய இயக்கிகளின் தொகுப்புகளில்
$\{ \lnot, \lor \}$, $\{ \lnot, \land \}$ மற்றும் $\{ \lnot,
\lif\}$ அடங்கும்---இவற்றை ஏதேனும் ஒரு அளவையடையுடன் சேர்த்தால்
வெளிப்பாட்டில் முழுமையான முதல்தர மொழி கிடைக்கும்.

ஷெஃபர் கோடு~$|$ (ஹென்றி ஷெஃபரின் பெயரால்) மற்றும் பியர்ஸின்
அம்பு~$\downarrow$, குவைனின் குத்துவாள் என்றும் அழைக்கப்படுவது, ஆகிய இரு
!!{operator}s உங்களுக்குத் தெரிந்திருக்கலாம். அவற்றின் வழக்கமான ``nand'' மற்றும்
``nor'' வாசிப்புகளில், இவ்வியக்கிகள் ஒவ்வொன்றும் தனியாகவே மெய்மைச் சார்பில்
முழுமையானவை.
\end{explain}

\end{document}
